% =============================================================================
%  Icare — manuel complet, rapport et documentation technique
%  Compilation : latexmk -pdf icare.tex   (ou voir build.ps1)
% =============================================================================
\documentclass[11pt,a4paper,oneside]{report}

% =============================================================================
%  Icare — préambule commun (paquets, couleurs, styles de code, macros)
% =============================================================================

% ---- Encodage & langue -------------------------------------------------------
\usepackage[T1]{fontenc}
\usepackage[utf8]{inputenc}
\usepackage[french]{babel}
\usepackage{lmodern}
\usepackage{textcomp}
\usepackage{microtype}

% Let TeX stretch a little before letting a line overflow into the margin.
\setlength{\emergencystretch}{2em}

% ---- Mise en page ------------------------------------------------------------
\usepackage[a4paper,margin=2.4cm,headheight=15pt]{geometry}
\usepackage{parskip}
\usepackage{fancyhdr}

% ---- Tableaux & listes -------------------------------------------------------
\usepackage{booktabs}
\usepackage{longtable}
\usepackage{array}
\usepackage{enumitem}
\usepackage{float}

% ---- Graphiques & couleurs ---------------------------------------------------
\usepackage{graphicx}
\usepackage[dvipsnames,table]{xcolor}
\usepackage{tikz}
\usetikzlibrary{arrows.meta,positioning,shapes.geometric,fit,backgrounds,calc}

% ---- Maths & symboles --------------------------------------------------------
\usepackage{amsmath,amssymb}

% ---- Code source -------------------------------------------------------------
\usepackage{listings}

% ---- Liens (en dernier) ------------------------------------------------------
\usepackage[hidelinks,unicode]{hyperref}
\hypersetup{
  pdftitle={Icare - manuel et documentation},
  pdfauthor={Icare contributors},
  pdfsubject={Controleur DMX web},
  pdfkeywords={DMX, Flask, lumiere, spectacle, musique}
}

% =============================================================================
%  Couleurs
% =============================================================================
\definecolor{icareblue}{HTML}{2D6CDF}
\definecolor{icaredark}{HTML}{1B3A6B}
\definecolor{icodebg}{HTML}{F4F6FA}
\definecolor{icodeln}{HTML}{C3CCDA}
\definecolor{icodecmt}{HTML}{6B7280}
\definecolor{icodestr}{HTML}{1F7A5A}
\definecolor{icodokw}{HTML}{8A2BE2}
\definecolor{icarenote}{HTML}{FDF6E3}
\definecolor{icarenborder}{HTML}{C9A227}

% =============================================================================
%  Style des listings
% =============================================================================
\lstdefinestyle{icare}{
  backgroundcolor=\color{icodebg},
  basicstyle=\ttfamily\small,
  keywordstyle=\color{icodokw}\bfseries,
  commentstyle=\color{icodecmt}\itshape,
  stringstyle=\color{icodestr},
  numberstyle=\tiny\color{icodeln},
  numbers=left,
  numbersep=9pt,
  frame=single,
  framerule=0.4pt,
  rulecolor=\color{icodeln},
  breaklines=true,
  breakatwhitespace=false,
  showstringspaces=false,
  tabsize=2,
  captionpos=b,
  columns=fullflexible,
  keepspaces=true,
  upquote=false,
}
\lstset{style=icare}

\lstdefinelanguage{json}{
  stringstyle=\color{icodestr},
  morestring=[b]",
  literate=
    *{0}{{{\color{icodokw}0}}}{1}
     {1}{{{\color{icodokw}1}}}{1}
     {2}{{{\color{icodokw}2}}}{1}
     {3}{{{\color{icodokw}3}}}{1}
     {4}{{{\color{icodokw}4}}}{1}
     {5}{{{\color{icodokw}5}}}{1}
     {6}{{{\color{icodokw}6}}}{1}
     {7}{{{\color{icodokw}7}}}{1}
     {8}{{{\color{icodokw}8}}}{1}
     {9}{{{\color{icodokw}9}}}{1},
}

% =============================================================================
%  En-têtes & pieds de page
% =============================================================================
\pagestyle{fancy}
\fancyhf{}
\fancyhead[L]{\textcolor{icaredark}{\textbf{Icare}}}
\fancyhead[R]{\textcolor{icodecmt}{\leftmark}}
\fancyfoot[C]{\textcolor{icodecmt}{\thepage}}
\renewcommand{\headrulewidth}{0.4pt}
\renewcommand{\footrulewidth}{0pt}
\fancypagestyle{plain}{%
  \fancyhf{}%
  \fancyfoot[C]{\textcolor{icodecmt}{\thepage}}%
  \renewcommand{\headrulewidth}{0pt}%
}

% =============================================================================
%  Macros
% =============================================================================
\newcommand{\code}[1]{\texttt{#1}}
\newcommand{\api}[1]{\texttt{#1}}
\newcommand{\token}[1]{\texttt{#1}}
\newcommand{\file}[1]{\texttt{#1}}
\newcommand{\Icare}{\textsc{Icare}}

% Encadré "à noter"
\newenvironment{note}{%
  \begin{center}%
  \begin{tabular}{@{}|p{0.92\linewidth}|@{}}
  \hline
  \rowcolor{icarenote}
  \textbf{À noter}\\[2pt]
}{%
  \\ \hline
  \end{tabular}%
  \end{center}%
}

% Encadré "attention"
\newenvironment{avertissement}{%
  \begin{center}%
  \begin{tabular}{@{}|p{0.92\linewidth}|@{}}
  \hline
  \rowcolor{red!10}
  \textbf{Attention}\\[2pt]
}{%
  \\ \hline
  \end{tabular}%
  \end{center}%
}

% Petite légende sous les figures
\newcommand{\legende}[1]{\par\smallskip{\small\itshape\textcolor{icodecmt}{#1}}\par}


\begin{document}

% -----------------------------------------------------------------------------
%  Page de titre
% -----------------------------------------------------------------------------
\begin{titlepage}
  \centering
  \vspace*{2.5cm}

  \begin{tikzpicture}
    \node[
      rounded corners=12pt,
      fill=icareblue!12,
      draw=icareblue,
      line width=1pt,
      inner sep=18pt,
    ] {\fontsize{58}{62}\selectfont\textbf{\textcolor{icaredark}{Icare}}};
  \end{tikzpicture}

  \vspace{1.2cm}
  {\Huge\bfseries Manuel et documentation}\\[0.4cm]
  {\Large\textcolor{icaredark}{Contrôleur DMX web pour groupes amateurs}}\\[0.3cm]
  {\large Version 0.2.0 --- octobre 2026}

  \vspace{1.8cm}
  \begin{tikzpicture}
    \node[draw=icodeln,fill=icodebg,rounded corners,inner sep=12pt,text width=0.78\linewidth,align=left]{
      \textbf{En une phrase.} \Icare{} transforme un adaptateur USB-DMX
      (compatible Enttec) en pupitre lumière pilotable depuis le navigateur :
      couleurs, intensité, effets, séquences et lecture sur scène en BPM.
    };
  \end{tikzpicture}

  \vfill
  {\large Auteur : \Icare{} contributors}\\[0.2cm]
  {\large Dépôt : \url{https://github.com/DlWAT/Icare}}\\[0.2cm]
  {\large Licence : MIT}
\end{titlepage}

% -----------------------------------------------------------------------------
%  Sommaire
% -----------------------------------------------------------------------------
\pagenumbering{roman}
\setcounter{tocdepth}{2}
\tableofcontents
\clearpage
\pagenumbering{arabic}

% -----------------------------------------------------------------------------
%  Corps du document
% -----------------------------------------------------------------------------
\part{Guide utilisateur}
% =============================================================================
%  PARTIE I — Guide utilisateur
% =============================================================================

\chapter{Présentation}
\label{ch:presentation}

\Icare{} est un petit contrôleur de lumière DMX qui se présente comme une
application web locale. Concrètement : un programme Python démarre un serveur
\emph{Flask} sur votre machine, et votre navigateur affiche l'interface. Un
adaptateur USB-DMX branché en USB assure la liaison physique avec les
projecteurs.

Le projet vise les \textbf{groupes amateurs} qui veulent se passer d'un vrai
pupitre et d'un éclairagiste dédié : on branche, on ouvre la page, et les
lumières répondent.

\section{Principes directeurs}

\begin{itemize}[leftmargin=1.4em]
  \item \textbf{100 \% local.} Rien ne part sur Internet ; le navigateur ne
        parle qu'à \code{127.0.0.1:5000}. Pas de compte, pas de cloud.
  \item \textbf{Une seule interface.} Le navigateur est le pupitre : trois
        onglets suffisent (\emph{Direct}, \emph{Séquence}, \emph{Live}).
  \item \textbf{Démarrage rapide.} Une commande pour installer, une commande pour
        lancer. Une version \code{.exe} autonome existe aussi (voir
        chapitre~\ref{ch:installation}).
  \item \textbf{Le spectaculaire accessible.} Effets intégrés (arc-en-ciel, feu,
        gyrophare, chenillard), mode \emph{audio-réactif} au microphone, et
        lecture calée au tempo avec le bouton \emph{TAP}.
\end{itemize}

\section{Vocabulaire}

Avant d'aller plus loin, fixons les mots employés dans ce document et dans
l'interface.

\begin{description}[leftmargin=1.6em,style=nextline]
  \item[Univers (DMX)] Ensemble des 512 canaux d'une ligne DMX. Chaque canal
        porte une valeur de 0 à 255.
  \item[Module / projecteur (\emph{fixture})] Un appareil branché sur la ligne,
        qui occupe un ou plusieurs canaux. Dans \Icare{}, les modules portent des
        noms comme \code{ParLED1} ou \code{Strobe1}.
  \item[Canal de base] Premier canal d'un projecteur. Si un PAR RGB démarre au
        canal 9, alors R, G et B occupent respectivement les canaux 9, 10 et 11.
  \item[Séquence] Une suite d'étapes (\emph{lignes}) décrivant l'état des
        projecteurs, enregistrée dans un fichier JSON.
  \item[Section] Étiquette donnée à un groupe de lignes (\code{Intro\_x4},
        \code{Couplet1}\dots) pour se repérer.
  \item[Beat / ligne] Une unité de temps de la séquence. Une ligne = un temps
        musical.
  \item[BPM] \emph{Beats per minute}. Cadence de lecture de la séquence Live.
\end{description}

\section{Les trois onglets}

La barre de navigation en haut de page donne accès aux trois espaces de travail,
résumés dans le tableau~\ref{tab:onglets}.

\begin{table}[H]
  \centering
  \begin{tabular}{@{}p{2.4cm}p{11.5cm}@{}}
    \toprule
    \textbf{Onglet} & \textbf{À quoi il sert} \\
    \midrule
    \textbf{Direct} & Pilotage immédiat : couleurs, intensité, pan/tilt,
      blackout, couleurs aléatoires, strobe, effets intégrés, audio-réactif et
      visualisation 3D. \\
    \addlinespace
    \textbf{Séquence} & Édition d'un show temps par temps : table de lignes par
      module, insertion d'effets \emph{Uniform}, \emph{Fondu} et \emph{Strobe},
      sauvegarde dans un fichier JSON. \\
    \addlinespace
    \textbf{Live} & Lecture d'une séquence sur scène : choix des modules, BPM
      réglable ou mesuré au \emph{TAP}, transport, barre de progression et vue 3D
      du faisceau. \\
    \bottomrule
  \end{tabular}
  \caption{Rôle des trois onglets.}
  \label{tab:onglets}
\end{table}

En complément, la \textbf{barre supérieure} (ou \emph{topbar}) affiche un
témoin d'état DMX et un bouton de bascule de thème clair/sombre.

\begin{figure}[H]
  \centering
  \begin{tikzpicture}[
      font=\small,
      box/.style={draw=icodeln,fill=icodebg,rounded corners,inner sep=8pt,align=center,minimum width=3.4cm},
      nav/.style={draw=icareblue,fill=icareblue!10,rounded corners,inner sep=6pt,align=center,minimum width=2.2cm},
      arrow/.style={-Latex,draw=icodecmt},
    ]
    \node[box,minimum width=13cm] (top) {Topbar : marque \textbf{Icare} \quad|\quad Accueil \quad Direct \quad Séquence \quad Live \quad|\quad témoin DMX \quad thème};
    \node[nav,below=1.1cm of top,yshift=0.1cm,xshift=-4.2cm] (d) {Direct};
    \node[nav,right=0.5cm of d] (s) {Séquence};
    \node[nav,right=0.5cm of s] (l) {Live};
    \node[box,below=0.9cm of s,minimum width=13cm] (content) {Contenu de l'onglet actif};
    \draw[arrow] (d) -- (content);
    \draw[arrow] (s) -- (content);
    \draw[arrow] (l) -- (content);
  \end{tikzpicture}
  \caption{Organisation générale de l'interface.}
  \label{fig:ui}
\end{figure}

\chapter{Installation et démarrage}
\label{ch:installation}

\section{Prérequis}

\begin{itemize}[leftmargin=1.4em]
  \item Windows (les autres systèmes fonctionnent en développement, mais le
        paquetage fourni cible Windows).
  \item Python 3.10 ou plus récent.
  \item Un adaptateur USB-DMX compatible \emph{Enttec USB Pro} (ou un clone qui
        implémente le même protocole).
\end{itemize}

\section{Démarrage depuis les sources}

\begin{lstlisting}[language=bash,caption={Installation et lancement en développement.}]
python -m venv .venv
.\.venv\Scripts\Activate.ps1
pip install -r requirements.txt
python run_webapp.py
\end{lstlisting}

Le serveur écoute sur \code{http://127.0.0.1:5000}. Ouvrez cette adresse dans
votre navigateur. En mode \emph{exécutable} empaqueté, le navigateur s'ouvre
automatiquement au démarrage.

\section{Connecter l'adaptateur DMX}

\begin{enumerate}[leftmargin=1.6em]
  \item Branchez l'adaptateur USB-DMX sur le PC et reliez-le à la ligne DMX des
        projecteurs (XLR 5 broches ou 3 broches selon votre câblage).
  \item Dans l'onglet \textbf{Direct}, renseignez le \emph{canal de base} du
        premier projecteur (9 par défaut), puis cliquez sur
        \textbf{Connexion DMX}.
  \item Le témoin en haut à droite passe au vert (\emph{DMX OK}).
\end{enumerate}

Si aucun port n'est indiqué, \Icare{} choisit automatiquement le premier port
COM dont la description contient \og USB \fg{}. Vous pouvez forcer un port en
le saisissant dans le champ \emph{Port} (par exemple \code{COM3}).

\section{Construire un exécutable Windows}

Pour distribuer \Icare{} sur un PC sans Python, utilisez le script fourni :

\begin{lstlisting}[language=bash,caption={Construction de l'exécutable autonome.}]
.\build_exe.ps1
\end{lstlisting}

Cela produit \file{dist\textbackslash IcareWeb.exe} (environ 30 Mo). Au
lancement, l'exécutable démarre le serveur, ouvre le navigateur et crée ses
fichiers de travail à côté de lui.

\section{Où \Icare{} écrit-il ses fichiers ?}

\begin{table}[H]
  \centering
  \begin{tabular}{@{}llp{7.4cm}@{}}
    \toprule
    \textbf{Fichier / dossier} & \textbf{Type} & \textbf{Contenu} \\
    \midrule
    \file{Sequences/} & dossiers & Shows en JSON (fournis + créés). \\
    \file{Palettes/} & dossier & Palettes de couleurs nommées. \\
    \file{icare\_settings.json} & fichier & Préférences persistantes (thème, port,
      canal de base, dernière séquence, favoris). \\
    \file{icare.log} & fichier & Journal rotatif (512 Ko $\times$ 3). \\
    \bottomrule
  \end{tabular}
  \caption{Fichiers produits à l'exécution.}
  \label{tab:fichiers}
\end{table}

En développement, ces fichiers apparaissent à la racine du projet. En version
empaquetée, ils sont écrits à côté de l'exécutable, jamais dans l'archive
lecture seule.

\chapter{L'onglet Direct}
\label{ch:direct}

L'onglet \textbf{Direct} est le pupitre manuel : tout ce qu'on y fait est envoyé
immédiatement au contrôleur DMX.

\section{Connexion}

Le panneau \emph{Connexion} contient :

\begin{itemize}[leftmargin=1.4em]
  \item \textbf{Canal de base (R)} : premier canal couleur du projecteur
        (0 à 509).
  \item \textbf{Port (optionnel)} : nom du port COM ; vide = auto-détection.
  \item \textbf{Connexion DMX} : ouvre le port série.
  \item \textbf{Mettre à jour canal de base} : change le canal de base sans
        rouvrir le port.
\end{itemize}

\section{Couleur et intensité}

Les curseurs \textbf{Rouge}, \textbf{Vert} et \textbf{Bleu} vont de 0 à 255 ;
le curseur \textbf{Amplitude} règle la luminosité globale. Le sélecteur de
couleur (color picker) pilote directement les trois canaux RGB. Un aperçu
rectangulaire montre la couleur résultante.

Le bouton \textbf{+ Favori} enregistre la couleur courante dans la palette de
favoris ; un clic sur une pastille la rappelle. La croix \code{x} retire un
favori. La liste est persistée dans les réglages (32 favoris maximum).

\section{Orientation}

Les projecteurs du patch sont \textbf{fixes} : ils n'ont ni pan ni tilt (aucune
tête mobile). Leur direction se règle une fois, au montage, par le champ
\code{angle} de chaque appareil dans le patch. L'onglet Direct ne propose donc
plus de réglage de position ni d'orientation.

\section{Les modes}

\paragraph{Blackout / Stop.} \textbf{Blackout} éteint tout immédiatement
(couleurs à 0, amplitude à 0). \textbf{Stop} arrête le mode en cours.

\paragraph{Couleurs aléatoires.} Trois paramètres :
\emph{fréquence} (Hz), \emph{lissage} (0--100 \%) et \emph{luminosité}
(0--255). Le lissage interpole progressivement vers la couleur cible pour un
rendu plus doux.

\paragraph{Strobe (overlay).} Un strobe qui \emph{se superpose} au mode courant
sans l'interrompre. Réglages : fréquence (Hz), rapport cyclique (\%),
canal DMX optionnel et durée (s). Si un canal est précisé, c'est ce canal qui
clignote ; sinon l'image entière est alternée entre la couleur courante et le
noir.

\paragraph{Bibliothèque d'effets.} Quatre effets procéduraux :
\textbf{Rainbow}, \textbf{Fire}, \textbf{Police} et \textbf{Chase RGB}, avec
réglage de \emph{vitesse} et de \emph{luminosité}. Voir le
chapitre~\ref{ch:effets} pour le détail.

\paragraph{Audio-réactif (bêta).} Le microphone du navigateur pilote
l'intensité au rythme de la musique. On règle la \emph{sensibilité} et la
\emph{couleur} ; une barre indique le niveau détecté. Voir le
chapitre~\ref{ch:audio}.

\section{Visualisation 3D}

Le panneau \emph{Visualisation 3D} dessine un projecteur, son faisceau et un
faisceau spot sur une scène. Couleur, intensité et orientation suivent les
curseurs (et les modes actifs). Pendant un strobe overlay, l'opacité du faisceau
clignote au même rythme.

\chapter{L'onglet Séquence}
\label{ch:sequence-ui}

L'onglet \textbf{Séquence} est l'éditeur de spectacle. Il manipule un fichier
JSON que l'on pourra lire ensuite dans l'onglet \emph{Live}.

\section{Fichiers}

Le panneau \emph{Fichiers} liste les séquences disponibles. Chaque entrée offre
un bouton \textbf{Charger} et une croix \code{x} de suppression. Le champ du
haut permet de nommer puis \textbf{Sauvegarder} le show courant.

\section{Barre d'édition}

\begin{table}[H]
  \centering
  \begin{tabular}{@{}lp{10.5cm}@{}}
    \toprule
    \textbf{Champ} & \textbf{Signification} \\
    \midrule
    \textbf{Module} & Projecteur ciblé par l'insertion d'effet. \\
    \textbf{\#Ligne} & Numéro de la ligne (temps) où insérer. \\
    \textbf{\#Tps} & Nombre de temps que dure l'effet (\code{n}). \\
    \textbf{Fréquence} & Fréquence du strobe, en Hz. \\
    Couleur 1 / Couleur 2 & Couleurs de départ et d'arrivée. \\
    \bottomrule
  \end{tabular}
  \caption{Contrôles de la barre d'édition.}
\end{table}

Trois boutons insèrent des effets :

\begin{itemize}[leftmargin=1.4em]
  \item \textbf{Uniform} : couleur fixe pendant \code{n} temps.
  \item \textbf{Fondu} : dégradé linéaire de la couleur 1 vers la couleur 2 sur
        \code{n} temps.
  \item \textbf{Strobe} : clignotement à la fréquence donnée pendant \code{n}
        temps.
\end{itemize}

Le bouton \textbf{+ Ligne} ajoute une ligne vierge, et \textbf{Effacer cellule}
vide la cellule ciblée.

\section{Édition directe du tableau}

Le tableau comporte une colonne par module. Les cellules de fonction sont
éditionnables directement : on peut y taper à la main un jeton du langage
décrit au chapitre~\ref{ch:dsl}. Le caractère \token{|} prolonge l'effet
commencé sur la ligne précédente.

\begin{note}
Une cellule vide signifie \og{}ne rien changer pour ce module à ce temps\fg{}.
C'est utile pour ne modifier qu'un seul projecteur à la fois.
\end{note}

\chapter{L'onglet Live}
\label{ch:live-ui}

L'onglet \textbf{Live} joue une séquence en conditions de spectacle.

\section{Charger et sélectionner}

\begin{enumerate}[leftmargin=1.6em]
  \item Choisissez une séquence dans la liste (\textbf{Charger}).
  \item Cochez les \textbf{modules} à piloter. Par défaut, tous les modules de la
        séquence sont pris en compte. Les boutons \emph{Tout sélectionner} et
        \emph{Tout désélectionner} font le tri.
\end{enumerate}

\section{Régler le tempo}

Deux méthodes :

\begin{itemize}[leftmargin=1.4em]
  \item Saisir un \textbf{BPM} (20 à 300) dans le champ numérique.
  \item Cliquer sur \textbf{TAP} en rythme : après deux appuis rapprochés,
        \Icare{} calcule le tempo moyen et l'applique. \textbf{Reset} efface
        l'historique des appuis.
\end{itemize}

\section{Transport}

\textbf{Play}, \textbf{Pause}, \textbf{Stop}, \textbf{Précédent} et
\textbf{Suivant} contrôlent la lecture. La barre de
progression, le nom de la section courante et le compteur \code{i / total}
donnent l'avancement. À la fin, la dernière image est maintenue puis la lecture
s'arrête.

\section{Vue 3D}

Le faisceau affiché prend la couleur envoyée au DMX en temps réel, ce qui permet
de vérifier ce qui part sur la ligne même sans regarder les projecteurs.

\chapter{Réglages et persistance}
\label{ch:reglages}

Peu de choses à régler, mais elles survivent aux redémarrages.

\begin{table}[H]
  \centering
  \begin{tabular}{@{}ll@{}}
    \toprule
    \textbf{Réglage} & \textbf{Effet} \\
    \midrule
    Thème & Clair ou sombre (bouton lune/soleil de la topbar). \\
    Dernier port COM & Pré-remplit le champ port. \\
    Canal de base & Restauré au démarrage. \\
    Dernière séquence & Mémorisée pour rechargement. \\
    Couleurs favorites & Palette affichée dans l'onglet Direct. \\
    \bottomrule
  \end{tabular}
  \caption{Réglages persistés.}
\end{table}

Tout est écrit dans \file{icare\_settings.json}. Le journal \file{icare.log}
enregistre les événements (connexion, erreurs série, etc.) et est automatiquement
limité en taille.

\begin{avertissement}
Le fichier de réglages est du JSON : s'il devient illisible (édition manuelle
ratée), \Icare{} le réinitialise silencieusement au démarrage suivant.
\end{avertissement}


\part{Comprendre \Icare}
% =============================================================================
%  PARTIE II — Comprendre Icare
% =============================================================================

\chapter{Le DMX512 en bref}
\label{ch:dmx}

Cette partie explique \emph{pourquoi} le code fait ce qu'il fait. On y revoit le
protocole DMX, la trame de l'adaptateur, la circulation des données et les
algorithmes des effets.

\section{Un univers de 512 canaux}

DMX512 est un protocole de scène très simple. Un \emph{univers} contient
\textbf{512 canaux}, chacun sur un octet, donc une valeur de \textbf{0 à 255}.
Onze bits sont envoyés par canal : une impulsion de \emph{start bit}, huit bits de
données, deux \emph{stop bits}. Le tout circule sur un bus RS-485 à
\textbf{250 000 bauds}.

Deux propriétés importantes :

\begin{itemize}[leftmargin=1.4em]
  \item \textbf{Aucun retour.} Le contrôleur parle, les projecteurs écoutent.
        On ne sait jamais si un projecteur a compris autre chose que ce qu'on
        voulait.
  \item \textbf{Image continue.} La même \emph{image} (les 512 valeurs) est
        renvoyée en boucle. Un projecteur qui cesse de recevoir se met en
        sécurité (souvent après quelques dixièmes de seconde).
\end{itemize}

Les valeurs n'ont pas de sens universel : c'est le projecteur qui décide, selon
son \emph{profil}. Un PAR RGB courant utilise trois canaux (\emph{intensité},
\emph{rouge}, \emph{vert}, \emph{bleu} sur certains modèles, parfois plus).

\section{Le plan de canaux d'\Icare{}}

Pour rester simple, \Icare{} pilote un unique projecteur RGB selon un plan de
canaux fixe, hérité de l'ancien prototype. Ces constantes vivent dans
\file{webapp/config.py}, classe \code{Channels}.

\begin{table}[H]
  \centering
  \begin{tabular}{@{}clp{8cm}@{}}
    \toprule
    \textbf{Canal} & \textbf{Constante} & \textbf{Rôle} \\
    \midrule
    0 & \code{MASTER} & Toujours à 255 pour «{}\/ouvrir\/»{} la sortie. \\
    1 & \code{AMPLITUDE} & Luminosité globale (0--255). \\
    3 & \code{PAN} & Orientation horizontale. \\
    5 & \code{TILT} & Orientation verticale. \\
    8 & \code{SHUTTER} & Toujours à 0 (pas d'effet de shutter). \\
    9 & \code{BASE\_COLOR\_DEFAULT} & R du projecteur, G en 10, B en 11. \\
    \bottomrule
  \end{tabular}
  \caption{Plan de canaux utilisé par le contrôleur.}
  \label{tab:channels}
\end{table}

Le \emph{canal de base} est modifiable : les canaux R, G, B sont alors
\code{base}, \code{base+1} et \code{base+2}. \Icare{} accepte une base de 0 à
509, afin qu'il reste toujours trois canaux disponibles.

\chapter{L'adaptateur USB-DMX}
\label{ch:enttec}

Le PC n'a pas de port DMX natif : on passe par un boîtier USB dont l'un des
plus répandus est l'\emph{Enttec USB Pro}. Le dialogue se fait à 250 000 bauds,
en 8 bits, sans parité, 2 bits de stop (8N2). \Icare{} occupe la classe
\code{DMX\_Icare} de \file{pydmxIcare.py}.

\section{La trame Enttec USB Pro}

Chaque envoi est encapsulé dans une petite trame :

\begin{center}
\begin{tikzpicture}[font=\ttfamily\small,every node/.style={draw,minimum height=0.9cm,minimum width=1.5cm,align=center}]
  \node[fill=icarenote] (a) {0x7E};
  \node[right=0pt of a,fill=icareblue!12] (b) {0x06};
  \node[right=0pt of b,fill=icareblue!12] (c) {LSB};
  \node[right=0pt of c,fill=icareblue!12] (d) {MSB};
  \node[right=0pt of d,fill=icarenote] (e) {0x00};
  \node[right=0pt of e,fill=green!12,minimum width=2.6cm] (f) {512 octets};
  \node[right=0pt of f,fill=icarenote] (g) {0xE7};
\end{tikzpicture}
\end{center}
\legende{Trame Enttec USB Pro : début, étiquette «{}Output Only Send DMX Packet
Request{}», longueur de charge utile en petit-boutiste, code de départ DMX,
données, fin.}

Détail des champs :

\begin{itemize}[leftmargin=1.4em]
  \item \code{0x7E} : octet de début de trame (\emph{START\_BYTE}).
  \item \code{0x06} : étiquette «{}envoi seul de paquet DMX{}»
        (\emph{LABEL\_DMX\_TX}).
  \item \textbf{LSB, MSB} : longueur de la charge utile = 513 octets
        (1 code de départ + 512 canaux), soit \code{0x01 0x02}.
  \item \code{0x00} : code de départ DMX standard.
  \item \textbf{512 octets} : les valeurs des canaux.
  \item \code{0xE7} : octet de fin de trame (\emph{END\_BYTE}).
\end{itemize}

Dans le code, l'en-tête est précalculé une fois pour toutes sous forme de
tableau NumPy, puis concaténé aux 512 canaux à chaque envoi :

\begin{lstlisting}[language=Python,caption={Construction et envoi d'une trame dans \file{pydmxIcare.py}.}]
payload_len = NUM_CHANNELS + 1  # 1 start code + 512 channels
self.start_byte = np.array([
    START_BYTE, LABEL_DMX_TX,
    payload_len & 0xFF, (payload_len >> 8) & 0xFF,
    0x00,  # DMX start code
], np.uint8)
...
frame = np.concatenate((self.start_byte, self.data, self.end_byte)).tobytes()
self.com.write(frame)
\end{lstlisting}

\section{Détection du port}

Si aucun port n'est fourni, \Icare{} parcourt les ports série et retient le
premier dont la description contient \og USB \fg{} :

\begin{lstlisting}[language=Python,caption={Auto-détection du port série.}]
for port_no, description, _address in serial.tools.list_ports.comports():
    if "USB" in description:
        return str(port_no)
return None
\end{lstlisting}

En cas d'échec d'écriture, le port est refermé proprement plutôt que de laisser
la boucle échouer en rafale.

\chapter{La chaîne logicielle, de bout en bout}
\label{ch:chaine}

Voici comment un mouvement de curseur finit en photons.

\begin{figure}[H]
  \centering
  \begin{tikzpicture}[
      node distance=8mm,
      font=\small,
      every node/.style={draw=icodeln,fill=icodebg,rounded corners,align=center,inner sep=6pt},
      flow/.style={-Latex,draw=icareblue,thick},
    ]
    \node (browser) {Navigateur\\(curseur, effets, audio)};
    \node[right=1.5cm of browser] (flask) {Flask\\API REST};
    \node[right=1.5cm of flask] (dmx) {\code{DMXController}\\thread de sortie};
    \node[below=1.1cm of dmx] (serial) {\code{pydmxIcare}\\port série};
    \node[left=1.5cm of serial] (adapter) {Adaptateur\\USB-DMX};
    \node[left=1.5cm of adapter] (fixtures) {Projecteurs};
    \draw[flow] (browser) -- node[above,draw=none,fill=none,font=\scriptsize]{HTTP JSON} (flask);
    \draw[flow] (flask) -- node[above,draw=none,fill=none,font=\scriptsize]{appel Python} (dmx);
    \draw[flow] (dmx) -- node[right,draw=none,fill=none,font=\scriptsize]{trame} (serial);
    \draw[flow] (serial) -- node[above,draw=none,fill=none,font=\scriptsize]{USB} (adapter);
    \draw[flow] (adapter) -- node[above,draw=none,fill=none,font=\scriptsize]{DMX512} (fixtures);
    \draw[flow,dashed] (fixtures.south) to[out=-90,in=-90,distance=3.2cm] node[below,draw=none,fill=none,font=\scriptsize]{lumières (aucun retour)} (browser.south);
  \end{tikzpicture}
  \caption{Chaîne complète, du navigateur aux projecteurs.}
  \label{fig:chain}
\end{figure}

Les étapes :

\begin{enumerate}[leftmargin=1.6em]
  \item Le navigateur (\file{app.js}) envoie par exemple
        \api{POST /api/experiment/set-values} avec \code{\{r, g, b, amp, pan, tilt\}}.
  \item La route correspondante (\file{webapp/api\_experiment.py}) appelle
        \code{dmx.set\_values(...)}.
  \item \code{DMXController} borne les valeurs (0--255) et, selon le mode actif,
        écrit la \emph{frame} courante.
  \item La boucle de fond (thread \code{dmx-out}) rafraîchit la sortie en
        continu (environ 20 Hz en mode stable).
  \item \code{pydmxIcare} assemble la trame Enttec et l'écrit sur le port série.
  \item L'adaptateur traduit cela en signal DMX512 sur la ligne.
\end{enumerate}

\section{Pourquoi un thread ?}

L'écriture série et les effets sont des tâches \emph{continues} qui ne doivent
pas bloquer les requêtes HTTP. \Icare{} isole donc la sortie DMX dans un thread
démon, et le lecteur de séquence dans un second thread. Les routes HTTP ne font
que modifier l'état partagé (couleurs, paramètres de mode) ; le thread s'occupe
de le refléter sur le fil.

\section{L'état partagé}

Le contrôleur maintient :

\begin{itemize}[leftmargin=1.6em,itemsep=2pt,topsep=3pt]
  \item \code{r, g, b, amp, pan, tilt} : les valeurs de sortie ;
  \item \code{enabled, strobe\_enabled, random\_enabled, effect\_enabled,
        audio\_enabled} : un seul mode actif à la fois ;
  \item \code{overlay} : le strobe superposé, indépendant du mode ;
  \item \code{random\_params, effect\_params, base\_color\_channel}.
\end{itemize}

\chapter{Les effets procéduraux}
\label{ch:effets}

Quatre effets sont générés côté serveur, dans la boucle de sortie. Le temps
écoulé depuis le démarrage, multiplié par la \emph{vitesse}, sert de paramètre
d'animation.

\begin{table}[H]
  \centering
  \begin{tabular}{@{}lp{10.5cm}@{}}
    \toprule
    \textbf{Effet} & \textbf{Principe} \\
    \midrule
    \code{rainbow} & Teinte qui tourne : $h = (60\,t) \bmod 360$ degrés, saturation
      et valeur maximales (conversion HSV$\rightarrow$RGB). \\
    \addlinespace
    \code{fire} & Orange/rouge vacillant : teinte autour de 18\textdegree, valeur
      aléatoire entre 0,55 et 1,0 pour l'effet de flamme. \\
    \addlinespace
    \code{police} & Alternance rouge / bleu à environ 4 Hz
      ($\lfloor 4t \rfloor \bmod 2$). \\
    \addlinespace
    \code{chase} & Chenillard qui parcourt R, G, B à environ 3 pas par seconde. \\
    \bottomrule
  \end{tabular}
  \caption{Effets procéduraux intégrés.}
\end{table}

La conversion HSV vers RGB est implémentée à la main, sans dépendance externe :

\begin{lstlisting}[language=Python,caption={Le cœur de \code{\_hsv\_to\_rgb}.}]
c = v * s
x = c * (1 - abs(((h / 60) % 2) - 1))
m = v - c
if   h <  60: r1, g1, b1 = c, x, 0
elif h < 120: r1, g1, b1 = x, c, 0
...
return int((r1 + m) * 255), int((g1 + m) * 255), int((b1 + m) * 255)
\end{lstlisting}

\section{Le strobe en superposition}

Le strobe d'overlay est à part : il ne remplace pas le mode courant, il passe
par-dessus. À chaque image, on évalue une phase
$\phi = (t \cdot f) \bmod 1$ et on éteint la sortie quand $\phi$ dépasse le
rapport cyclique :

\begin{equation*}
\text{sortie} =
\begin{cases}
\text{couleur courante} & \text{si } \phi < \text{duty}, \\
\text{noir (ou 0 sur le canal choisi)} & \text{sinon.}
\end{cases}
\end{equation*}

Si un canal DMX est précisé, seul ce canal clignote ; sinon toute l'image
alterne. L'overlay s'arrête tout seul à la fin de la durée demandée.

\chapter{Le mode audio-réactif}
\label{ch:audio}

Ce mode transforme le microphone du navigateur en capteur de rythme. Toute
l'analyse se fait côté client ; le serveur ne reçoit que le résultat.

\begin{figure}[H]
  \centering
  \begin{tikzpicture}[
      node distance=9mm,
      font=\scriptsize,
      every node/.style={draw=icodeln,fill=icodebg,rounded corners,align=center,inner sep=5pt},
      flow/.style={-Latex,draw=icareblue,thick},
    ]
    \node (mic) {Micro\\\code{getUserMedia}};
    \node[right=0.35cm of mic] (an) {\code{Analyser}\\\code{fftSize=1024}};
    \node[right=0.35cm of an] (rms) {RMS\\\code{getByteTimeDomainData}};
    \node[right=0.35cm of rms] (lvl) {niveau $\times$ gain\\borne à $[0,1]$};
    \node[right=0.35cm of lvl] (api) {\api{POST}\\\code{/audio/push}};
    \node[below=1.2cm of api] (dmx) {\code{amp = niveau $\times$ 255}};
    \draw[flow] (mic) -- (an);
    \draw[flow] (an) -- (rms);
    \draw[flow] (rms) -- (lvl);
    \draw[flow] (lvl) -- (api);
    \draw[flow] (api) -- (dmx);
  \end{tikzpicture}
  \caption{Chaîne audio-réactive.}
  \label{fig:audio}
\end{figure}

\paragraph{Du son au niveau.} On échantillonne la forme d'onde, on recentre les
octets autour de 128, puis on calcule la valeur quadratique moyenne (RMS) :

\begin{equation*}
\mathrm{RMS} = \sqrt{\frac{1}{N}\sum_{i=0}^{N-1} \left(\frac{d_i - 128}{128}\right)^2 },
\qquad
\text{niveau} = \min\!\big(1,\ \mathrm{RMS} \times \text{gain}\big).
\end{equation*}

Le curseur \emph{Sensibilité} règle le \code{gain} entre 0,5 et 6.

\paragraph{Fréquence d'envoi.} La boucle d'animation tourne à la fréquence de
l'écran, mais l'envoi réseau est limité à environ 30 Hz pour ne pas saturer le
serveur. La couleur choisie accompagne le niveau.

\paragraph{Côté serveur.} \code{audio\_push(niveau, couleur)} met à jour
l'amplitude à \code{niveau $\times$ 255}, applique éventuellement la couleur, et
renvoie l'image. Le mode ne génère rien tout seul : sans échantillons, l'état
ne change pas.


\part{Documentation technique}
% =============================================================================
%  PARTIE III — Documentation technique
% =============================================================================

\chapter{Architecture du projet}
\label{ch:architecture}

\section{Arborescence}

\begin{lstlisting}[caption={Arborescence du dépôt.},numbers=none]
Icare/
|-- pydmxIcare.py          # bas niveau : trame Enttec + port serie
|-- pydmxlight.py          # descripteurs de projecteurs (Par LED, Strobe)
|-- run_webapp.py          # point d'entree (python run_webapp.py)
|-- webapp/
|   |-- __init__.py        # fabrique create_app()
|   |-- config.py          # constantes, UserSettings, journalisation
|   |-- services.py        # singletons (settings, dmx)
|   |-- utils.py           # helpers partages (noms de fichiers surs)
|   |-- dmx_controller.py  # machine a etats + thread de sortie DMX
|   |-- live_player.py     # lecteur de sequence + tap-tempo
|   |-- dsl.py             # analyseur du langage de sequences
|   |-- schemas.py         # validation / migration des fichiers JSON
|   |-- pages.py           # routes HTML
|   |-- api_experiment.py  # API onglet Direct
|   |-- api_sequence.py    # API editeur de sequence
|   |-- api_live.py        # API lecture Live
|   |-- api_health.py      # /api/health
|   |-- api_palettes.py    # favoris + palettes nommees
|   |-- api_settings.py    # preferences
|   |-- templates/         # gabarits Jinja2
|   `-- static/            # app.css, app.js, vendor/three.min.js
|-- Sequences/             # shows d'exemple
|-- tests/                 # suite pytest
|-- docs/                  # cette documentation (LaTeX)
|-- pyproject.toml         # build, ruff, pytest
`-- build_exe.ps1          # paquetage PyInstaller
\end{lstlisting}

\section{La fabrique d'application}

\code{create\_app()} construit et configure l'application Flask. Elle choisit
le dossier des gabarits et des fichiers statiques selon qu'on tourne depuis les
sources ou depuis un exécutable empaqueté (\code{sys.\_MEIPASS}), puis enregistre
les \emph{blueprints} un par un :

\begin{lstlisting}[language=Python,caption={Enregistrement des blueprints (\file{webapp/\_\_init\_\_.py}).}]
from .pages import pages_bp
from .api_experiment import api_experiment_bp
...
app.register_blueprint(pages_bp)
app.register_blueprint(api_experiment_bp)
app.register_blueprint(api_sequence_bp)
app.register_blueprint(api_live_bp)
app.register_blueprint(api_health_bp)
app.register_blueprint(api_palettes_bp)
app.register_blueprint(api_settings_bp)
\end{lstlisting}

Chaque famille de routes vit dans son propre module : on retrouve donc
naturellement une route \api{/api/experiment/...} dans
\file{api\_experiment.py}, et ainsi de suite.

\section{Les singletons de service}

Pour éviter les imports circulaires et garantir une seule instance partagée,
\file{webapp/services.py} expose :

\begin{itemize}[leftmargin=1.4em]
  \item \code{settings} : les préférences chargées depuis le disque ;
  \item \code{dmx} : l'unique \code{DMXController} du processus ;
  \item \code{SEQ\_DIR} / \code{PALETTE\_DIR} : les chemins des dossiers de
        données.
\end{itemize}

\begin{note}
Il n'existe qu'\textbf{un} contrôleur DMX et \textbf{un} lecteur Live pour tout
le serveur. C'est voulu : un seul fil pilote la ligne DMX.
\end{note}

\section{Configuration}

\file{webapp/config.py} distingue deux choses :

\begin{itemize}[leftmargin=1.4em]
  \item les \textbf{constantes} : plan de canaux (\code{Channels}) et cadences
        de boucle (\code{Timings}) ;
  \item l'\textbf{état persistant} : la dataclass \code{UserSettings}.
\end{itemize}

\code{UserSettings} sait se charger, se sauvegarder et mettre à jour des champs.
Au chargement, seules les clés connues sont conservées, ce qui rend le fichier
tolérant aux évolutions de schéma. La journalisation écrit à la fois sur la
console et dans un fichier rotatif.

\begin{table}[H]
  \centering
  \begin{tabular}{@{}lll@{}}
    \toprule
    \textbf{Cadence} & \textbf{Valeur} & \textbf{Rôle} \\
    \midrule
    \code{STEADY\_LOOP} & 0,05 s & Sortie stable, environ 20 Hz. \\
    \code{RANDOM\_LOOP} & 0,025 s & Environ 40 Hz pour des fondus fluides. \\
    \code{EFFECT\_LOOP} & 0,025 s & Environ 40 Hz pour les effets. \\
    \code{IDLE\_LOOP} & 0,10 s & Attente quand aucun mode n'est actif. \\
    \bottomrule
  \end{tabular}
  \caption{Cadences de la boucle DMX (\code{Timings}).}
\end{table}

\section{Le front-end}

\file{webapp/static/app.js} est un unique script sans dépendance (hormis
Three.js, fourni en local dans \file{vendor/}). Au chargement, il lit
\code{document.body.dataset.page} et n'active que le module correspondant
(\code{bindExperiment}, \code{bindSequence} ou \code{bindLive}). Deux sondages
réguliers maintiennent l'interface synchronisée avec l'état serveur :

\begin{itemize}[leftmargin=1.4em]
  \item \api{GET /api/experiment/status} toutes les 200 ms (onglet Direct) ;
  \item \api{GET /api/live/status} toutes les 250 ms (onglet Live).
\end{itemize}

\chapter{Le contrôleur DMX}
\label{ch:controleur}

\code{DMXController} (\file{webapp/dmx\_controller.py}) est le cœur temps réel.

\section{Des modes exclusifs, un overlay}

Un seul mode peut être actif à la fois. Le strobe d'overlay, lui, est
indépendant et se superpose.

\begin{center}
\begin{tikzpicture}[
    font=\small,
    mode/.style={draw=icareblue,fill=icareblue!10,rounded corners,inner sep=7pt},
    arrow/.style={-Latex,draw=icodecmt},
  ]
  \node[mode] (idle) {inactif};
  \node[mode,right=1.6cm of idle] (steady) {steady};
  \node[mode,above=0.8cm of steady] (strobe) {strobe};
  \node[mode,below=0.8cm of steady] (random) {random};
  \node[mode,below=0.8cm of random] (effect) {effect};
  \node[mode,below=0.8cm of effect] (audio) {audio};
  \draw[arrow] (idle) -- (steady);
  \draw[arrow] (steady) -- (strobe);
  \draw[arrow] (strobe) -- (random);
  \draw[arrow] (random) -- (effect);
  \draw[arrow] (effect) -- (audio);
  \draw[arrow,dashed,bend left=35] (strobe.north) to node[above,draw=none,fill=none,font=\scriptsize]{switch} (idle.north);
\end{tikzpicture}
\end{center}
\legende{\code{\_switch\_mode()} réinitialise tous les drapeaux puis active le
mode voulu : les modes sont mutuellement exclusifs.}

\section{Le thread de sortie}

Le thread \code{dmx-out} boucle sans fin et exécute une étape selon le mode
actif :

\begin{lstlisting}[language=Python,caption={Boucle principale du contrôleur (simplifiée).}]
while not self._stop_event.is_set():
    try:
        if self.enabled:          self._send_with_overlay(); time.sleep(Timings.STEADY_LOOP)
        elif self.strobe_enabled: self._step_strobe()
        elif self.random_enabled: self._step_random()
        elif self.effect_enabled: self._step_effect()
        elif self.audio_enabled:  time.sleep(Timings.IDLE_LOOP)
        else:                     time.sleep(Timings.IDLE_LOOP)
    except Exception:
        logger.exception("DMX loop error; pausing 0.5 s")
        time.sleep(0.5)
\end{lstlisting}

Chaque \emph{frame} passe par \code{\_send\_with\_overlay()}, qui applique
d'abord l'overlay éventuel puis envoie l'image. Les écritures série sont
encapsulées dans des \code{try/except} : une erreur transitoire est journalisée
et la boucle repart après une courte pause, au lieu de tuer le thread.

\section{Points de robustesse}

\begin{itemize}[leftmargin=1.4em]
  \item Toutes les entrées sont \textbf{bornées} : les composantes RGB,
        l'amplitude, le pan et le tilt sont contraints à 0--255.
  \item Le canal de base est validé (0--509) et refusé sinon.
  \item \code{pydmxIcare.send\_data()} ne fait rien si le port est fermé : on
        peut lancer l'application sans matériel, sans plantage.
  \item En cas d'erreur d'écriture série, le port est refermé pour permettre une
        reconnexion propre.
\end{itemize}

\chapter{Le lecteur de séquence (LivePlayer)}
\label{ch:liveplayer}

\code{LivePlayer} (\file{webapp/live\_player.py}) parcourt les lignes d'une
séquence et traduit chaque cellule en couleurs.

\section{Chargement}

\code{load(raw)} normalise la séquence (voir chapitre~\ref{ch:schemas}) puis
remet l'index et l'état de lecture à zéro. Aucune lecture ne démarre tant que
\code{start()} n'est pas appelé.

\section{Transport}

\begin{table}[H]
  \centering
  \begin{tabular}{@{}lp{9.4cm}@{}}
    \toprule
    \textbf{Méthode} & \textbf{Effet} \\
    \midrule
    \code{start(bpm)} & Démarre la lecture (BPM optionnel). \\
    \code{pause()} & Bascule l'état \emph{pause}. \\
    \code{stop()} & Arrête et revient au début. \\
    \code{next()} / \code{prev()} & Avance ou recule d'une ligne. \\
    \code{toggle\_module(name)} & (Dé)sélectionne un module. \\
    \code{select\_all()} / \code{unselect\_all()} & Sélection globale. \\
    \bottomrule
  \end{tabular}
  \caption{API de transport du lecteur.}
\end{table}

\section{Boucle de lecture}

Le thread \code{live-player} applique une ligne puis attend la durée d'un temps :
\code{60 / BPM} secondes. À la fin de la séquence, la dernière image est
maintenue avant l'arrêt.

\begin{lstlisting}[language=Python,caption={Cœur de la boucle de lecture.}]
self.index = min(self.index + 1, total)
line = self.data["data"].get(str(self.index), {})
self._apply_step(line.get("function", []))
time.sleep(60.0 / max(1, self.bpm))
if self.index >= total:
    self.playing = False
\end{lstlisting}

\section{Sélection de modules et visualisation}

Si l'utilisateur a coché des modules, seuls ceux-là sont appliqués ; sinon tous
le sont. Pour la vue 3D, le lecteur retient la couleur du \emph{premier} module
ciblé et la pousse au contrôleur DMX, qui est mono-projecteur.

\section{Tap-tempo}

Le tempo mesuré vit dans \file{live\_player.py}. On conserve les horodatages
récents (fenêtre de 3 secondes), on calcule les intervalles, on en fait la
moyenne, et on borne le résultat entre 20 et 300 BPM :

\begin{lstlisting}[language=Python,caption={Calcul du tempo à partir des appuis.}]
intervals = [b - a for a, b in zip(self._taps, self._taps[1:], strict=False)]
avg = sum(intervals) / len(intervals)
bpm = int(round(60 / avg))
bpm = max(20, min(300, bpm))
\end{lstlisting}

\chapter{Le langage des séquences (DSL)}
\label{ch:dsl}

Chaque cellule d'une ligne contient une chaîne qui décrit l'action d'\textbf{un}
projecteur pour ce temps. \file{webapp/dsl.py} la transforme en objet typé.

\section{Grammaire}

\begin{table}[H]
  \centering
  \begin{tabular}{@{}ll@{}}
    \toprule
    \textbf{Jeton} & \textbf{Signification} \\
    \midrule
    \token{""} & Ne rien changer (garder l'état courant). \\
    \token{|} & Continuer l'effet multi-temps précédent. \\
    \token{uni\_n\_R\_G\_B} & Couleur fixe pendant \code{n} temps. \\
    \token{fond\_n\_R1\_G1\_B1\_R2\_G2\_B2} & Fondu linéaire de $c_1$ vers $c_2$. \\
    \token{str\_n\_R\_G\_B\_freq} & Strobe à \code{freq} Hz pendant \code{n} temps. \\
    \bottomrule
  \end{tabular}
  \caption{Jetons du langage de séquences.}
  \label{tab:dsl}
\end{table}

La fréquence du strobe est optionnelle : sans elle, la valeur par défaut est
60 Hz. Les longueurs \code{n} inférieures à 1 sont ramenées à 1, et toute
composante hors de 0--255 lève une \code{DSLError}.

\section{Structure du code}

Le module expose des dataclasses gelées :

\begin{itemize}[leftmargin=1.4em]
  \item \code{UniformEffect(length, color)} ;
  \item \code{FadeEffect(length, color\_from, color\_to)} ;
  \item \code{StrobeEffect(length, color, freq)} ;
  \item \code{ContinueEffect()} pour le jeton \token{|}.
\end{itemize}

La fonction \code{parse(token)} renvoie l'objet correspondant, ou \code{None}
pour une cellule vide. Trois \emph{encodeurs} (\code{encode\_uniform},
\code{encode\_fade}, \code{encode\_strobe}) font l'opération inverse, utilisés
notamment par la migration d'anciens fichiers.

\section{Le jeton de continuation}

Un effet de \code{n} temps s'écrit sur une première ligne, puis \code{n-1} fois
\token{|}. Le lecteur conserve, pour chaque module, un petit état
(mode, longueur, pas courant, couleurs, fréquence). À chaque \token{|}, il
avance d'un pas :

\begin{itemize}[leftmargin=1.4em]
  \item \textbf{fondu} : interpolation linéaire
        $c = c_1 + \frac{\text{pas}}{n-1}\,(c_2 - c_1)$ ;
  \item \textbf{strobe} : alternance couleur / noir à chaque pas ;
  \item \textbf{couleur fixe} : la couleur reste identique.
\end{itemize}

\begin{lstlisting}[caption={Exemple : un fondu de 4 temps du rouge au bleu.},numbers=none]
fond_4_255_0_0_0_0_255   |   |   |
\end{lstlisting}

\chapter{Schémas de fichiers et migration}
\label{ch:schemas}

Deux formats de séquence sont acceptés. \file{webapp/schemas.py} les unifie en
un seul : le format \emph{moderne}.

\section{Format moderne}

\begin{lstlisting}[language=json,caption={Séquence au format moderne.}]
{
  "Header": {
    "Structure": ["Intro_x4", "Intro_x4", "Couplet1_A_x2"],
    "Tempo": "60",
    "Modules": ["ParLED1", "ParLED2", "Strobe1"],
    "Line": ["1", "2", "3"]
  },
  "data": {
    "1": { "section": "Intro_x4", "function": ["", "uni_4_78_243_86", ""] },
    "2": { "section": "Intro_x4", "function": ["", "|", "str_3_132_214_107_63"] }
  }
}
\end{lstlisting}

\section{Format historique}

Les anciens fichiers (comme les premières versions de
\file{Doctor\_doctor.json}) ont une autre forme : une clé \code{headers} avec une
\code{structure}, et des lignes de quatre nombres
\code{[durée\_beats, R, G, B]}.

\begin{lstlisting}[language=json,caption={Ancien format (migré automatiquement).},numbers=none]
{
  "headers": { "structure": ["intro", "Couplet1", "Refrain1"] },
  "data": { "intro": [[2, 125, 20, 35], [2, 20, 53, 128]] }
}
\end{lstlisting}

\section{Normalisation}

\code{normalize\_sequence(raw)} renvoie toujours le format moderne :

\begin{itemize}[leftmargin=1.4em]
  \item les modules manquants sur une ligne sont complétés par des cellules
        vides, ou tronqués s'il y en a trop ;
  \item les clés de lignes sont converties en chaînes \code{"1"}, \code{"2"},
        \dots{} et triées numériquement ;
  \item en format historique, chaque durée est multipliée par deux (une durée de
        0,5 donne un temps, une durée de 2 en donne quatre) ; la première ligne
        reçoit un jeton \code{uni\_...}, les suivantes \token{|} ;
  \item un module factice \code{ParLED1} est utilisé, l'ancien format ne
        distinguant pas les projecteurs.
\end{itemize}

Si \code{raw} n'a ni \code{Header} ni \code{headers}, une \code{ValueError} est
levée : c'est ce qui alimente les réponses d'erreur 422 de l'API.

\begin{note}
À la sauvegarde, l'API valide en faisant un aller-retour de normalisation avant
d'écrire le fichier : un contenu invalide est donc refusé au lieu d'être
persisté.
\end{note}


\part{API, tests et déploiement}
% =============================================================================
%  PARTIE IV — API, tests et déploiement
% =============================================================================

\chapter{API REST}
\label{ch:api}

L'interface et le serveur communiquent exclusivement en JSON.

\section{Conventions}

\begin{itemize}[leftmargin=1.4em]
  \item \textbf{Succès} : \code{\{"ok": true, ...\}}, code HTTP \code{200}.
  \item \textbf{Échec} : \code{\{"ok": false, "error": "..."\}}, avec un code
        explicite :
        \begin{itemize}
          \item \code{400} : requête mal formée ou paramètre invalide ;
          \item \code{404} : fichier introuvable ;
          \item \code{422} : contenu illisible ou schéma inconnu ;
          \item \code{500} : erreur d'entrée/sortie côté serveur ;
          \item \code{503} : connexion DMX impossible.
        \end{itemize}
  \item Les noms de fichiers envoyés sont nettoyés et confinés dans
        \file{Sequences/} ou \file{Palettes/} : aucune traversée de répertoire
        n'est possible.
\end{itemize}

\section{Pages HTML}

\begin{table}[H]
  \centering
  \begin{tabular}{@{}lll@{}}
    \toprule
    \textbf{Méthode} & \textbf{Route} & \textbf{Page} \\
    \midrule
    GET & \api{/} & Accueil \\
    GET & \api{/experiment} & Onglet Direct \\
    GET & \api{/sequence} & Onglet Séquence \\
    GET & \api{/live} & Onglet Live \\
    \bottomrule
  \end{tabular}
  \caption{Routes HTML.}
\end{table}

\section{Onglet Direct}

\begin{longtable}{@{}llp{6.6cm}@{}}
  \toprule
  \textbf{Méthode} & \textbf{Route} & \textbf{Rôle} \\
  \midrule
  \endfirsthead
  \toprule
  \textbf{Méthode} & \textbf{Route} & \textbf{Rôle} \\
  \midrule
  \endhead
  POST & \api{/api/experiment/connect} & Ouvre le port (canal de base + port). \\
  POST & \api{/api/experiment/base-channel} & Change le canal de base. \\
  POST & \api{/api/experiment/set-values} & RGB, amplitude, pan, tilt. \\
  POST & \api{/api/experiment/start} & Démarre la sortie stable. \\
  POST & \api{/api/experiment/stop} & Arrête le mode courant. \\
  POST & \api{/api/experiment/blackout} & Éteint tout. \\
  POST & \api{/api/experiment/strobe} & Strobe plein cadre (freq, shuffle). \\
  POST & \api{/api/experiment/random} & Configure et démarre l'aléatoire. \\
  POST & \api{/api/experiment/mode/random-config} & Aléatoire : paramètres seuls. \\
  POST & \api{/api/experiment/random-start} & Démarre l'aléatoire. \\
  POST & \api{/api/experiment/random-stop} & Arrête l'aléatoire. \\
  POST & \api{/api/experiment/mode/strobe-overlay} & Strobe superposé. \\
  POST & \api{/api/experiment/effect/start} & Lance un effet (rainbow, \dots). \\
  POST & \api{/api/experiment/effect/stop} & Arrête l'effet. \\
  POST & \api{/api/experiment/audio/start} & Active le mode audio. \\
  POST & \api{/api/experiment/audio/stop} & Désactive le mode audio. \\
  POST & \api{/api/experiment/audio/push} & Pousse un échantillon (niveau, couleur). \\
  GET & \api{/api/experiment/status} & État complet (polling 200 ms). \\
  \bottomrule
  \caption{Points d'entrée de l'onglet Direct.}
\end{longtable}

Exemple de corps pour \api{set-values} :

\begin{lstlisting}[language=json]
{ "r": 255, "g": 128, "b": 0, "amp": 200, "pan": 90, "tilt": 45 }
\end{lstlisting}

\section{Séquence, Live, palettes, réglages}

\begin{longtable}{@{}llp{6.6cm}@{}}
  \toprule
  \textbf{Méthode} & \textbf{Route} & \textbf{Rôle} \\
  \midrule
  \endfirsthead
  \toprule
  \textbf{Méthode} & \textbf{Route} & \textbf{Rôle} \\
  \midrule
  \endhead
  \multicolumn{3}{@{}l}{\textbf{Séquence}}\\
  GET & \api{/api/sequence/list} & Liste les fichiers JSON. \\
  GET & \api{/api/sequence/load?name=} & Charge et normalise. \\
  POST & \api{/api/sequence/save} & Sauve (validation par aller-retour). \\
  POST & \api{/api/sequence/delete} & Supprime un fichier. \\
  \addlinespace
  \multicolumn{3}{@{}l}{\textbf{Live}}\\
  GET & \api{/api/live/list} & Liste les séquences lisibles. \\
  POST & \api{/api/live/load} & Charge une séquence dans le lecteur. \\
  POST & \api{/api/live/play} & Démarre (BPM optionnel). \\
  POST & \api{/api/live/pause} & Bascule pause. \\
  POST & \api{/api/live/stop} & Arrête et rembobine. \\
  POST & \api{/api/live/next} / \api{prev} & Ligne suivante / précédente. \\
  POST & \api{/api/live/toggle-module} & (Dé)sélectionne un module. \\
  POST & \api{/api/live/select-all} / \api{unselect-all} & Sélection globale. \\
  POST & \api{/api/live/tap} & Enregistre un appui de tempo. \\
  POST & \api{/api/live/tap-reset} & Efface l'historique des appuis. \\
  POST & \api{/api/live/bpm} & Fixe le BPM (20--300). \\
  GET & \api{/api/live/status} & Progression, section, BPM, couleur. \\
  \addlinespace
  \multicolumn{3}{@{}l}{\textbf{Palettes}}\\
  GET & \api{/api/palettes/favorites} & Liste des couleurs favorites. \\
  POST & \api{/api/palettes/favorites/add} & Ajoute un favori. \\
  POST & \api{/api/palettes/favorites/remove} & Retire un favori. \\
  GET & \api{/api/palettes/list} & Liste les palettes nommées. \\
  GET & \api{/api/palettes/load?name=} & Charge une palette. \\
  POST & \api{/api/palettes/save} & Enregistre une palette. \\
  POST & \api{/api/palettes/delete} & Supprime une palette. \\
  \addlinespace
  \multicolumn{3}{@{}l}{\textbf{Divers}}\\
  GET & \api{/api/settings} & Réglages courants. \\
  POST & \api{/api/settings} & Met à jour (thème, port, \dots). \\
  GET & \api{/api/health} & État du service et du DMX. \\
  \bottomrule
  \caption{Points d'entrée restants.}
\end{longtable}

\chapter{Tests, qualité et déploiement}
\label{ch:tests}

\section{Suite de tests}

Les tests utilisent \code{pytest} et ne nécessitent \textbf{aucun matériel} :
le port série est simplement fermé, donc les envois sont sans effet.

\begin{lstlisting}[language=bash,caption={Lancer la suite.}]
pip install -r requirements-dev.txt
pytest
\end{lstlisting}

\begin{table}[H]
  \centering
  \begin{tabular}{@{}lp{9.5cm}c@{}}
    \toprule
    \textbf{Fichier} & \textbf{Ce qui est vérifié} & \textbf{Tests} \\
    \midrule
    \file{test\_api.py} & Routes HTTP, bornage, sécurité des noms, favoris, réglages. & 10 \\
    \file{test\_dmx\_controller.py} & Bornage RGB/amp/pan/tilt, canal de base, HSV, effets inconnus. & 7 \\
    \file{test\_dsl.py} & Analyse, erreurs, bornes, aller-retours des encodeurs. & 16 \\
    \file{test\_schemas.py} & Passe-plat moderne, remplissage, migration historique. & 4 \\
    \midrule
    \multicolumn{2}{r}{\textbf{Total}} & \textbf{37} \\
    \bottomrule
  \end{tabular}
  \caption{Couverture fonctionnelle des tests.}
\end{table}

\section{Qualité du code}

Le projet est configuré pour \code{ruff}, à la fois pour le lint et pour le
formatage :

\begin{lstlisting}[language=bash,caption={Vérifications habituelles.}]
ruff check .          # lint
ruff format .         # formatage
pytest                # tests
\end{lstlisting}

Les règles activées couvrent les erreurs (\code{E}, \code{F}), les avertissements
(\code{W}), le tri des imports (\code{I}), les pièges fréquents (\code{B}),
la modernisation des annotations (\code{UP}) et les simplifications
(\code{C4}). La longueur de ligne cible est de 100 caractères.

\section{Construire l'exécutable}

Le script \file{build\_exe.ps1} et le fichier \file{IcareWeb.spec} emballent
l'application et ses ressources (gabarits, statiques, séquences) dans un
exécutable unique.

\begin{lstlisting}[language=bash,caption={Construction.}]
.\build_exe.ps1     # produit dist\IcareWeb.exe
\end{lstlisting}

\section{Sources et archive}

Pour ne pas embarquer de fichiers inutiles, \file{.gitignore} exclut les
caches Python (\code{\_\_pycache\_\_}, \code{*.pyc}), les environnements
virtuels, les artefacts PyInstaller (\file{build/}, \file{dist/}, \file{*.spec})
et les fichiers produits à l'exécution (\file{icare.log},
\file{icare\_settings.json}, \file{Palettes/}).

\begin{note}
Le dossier \file{Palettes/} est ignoré par Git car il contient des données
locales, mais il peut être sauvegardé à part si vous voulez conserver vos
couleurs favorites.
\end{note}


\appendix
% =============================================================================
%  ANNEXES
% =============================================================================

\chapter{Aide-mémoire du DSL}
\label{annexe:dsl}

\begin{longtable}{@{}p{4.7cm}p{4.4cm}p{4.9cm}@{}}
  \toprule
  \textbf{Jeton} & \textbf{Exemple} & \textbf{Effet} \\
  \midrule
  \endfirsthead
  \toprule
  \textbf{Jeton} & \textbf{Exemple} & \textbf{Effet} \\
  \midrule
  \endhead
  \token{""} & \token{""} & Ne change rien ce temps-là. \\
  \token{|} & \token{|} & Continue l'effet multi-temps précédent. \\
  \token{uni\_n\_R\_G\_B} & \token{uni\_4\_255\_0\_0} &
    Rouge fixe pendant 4 temps. \\
  \token{fond\_n\_c1\_c2} & \token{fond\_8\_255\_0\_0\_0\_0\_255} &
    Fondu rouge $\rightarrow$ bleu sur 8 temps. \\
  \token{str\_n\_R\_G\_B\_f} & \token{str\_2\_255\_255\_0\_120} &
    Strobe jaune pendant 2 temps à 120 Hz. \\
  \token{str\_n\_R\_G\_B} & \token{str\_3\_0\_0\_255} &
    Strobe bleu, fréquence par défaut (60 Hz). \\
  \bottomrule
  \caption{Tous les jetons acceptés.}
\end{longtable}

\paragraph{Règles.} Une composante de couleur doit être comprise entre 0 et
255, sans quoi le jeton est rejeté. La longueur \code{n} est au minimum de 1.
Un préfixe inconnu, un mauvais nombre d'arguments ou un texte non numérique
produisent une \code{DSLError}, journalisée et ignorée pendant la lecture Live.

\paragraph{Fondu sur plusieurs lignes.} On écrit le jeton complet sur la
première ligne, puis un \token{|} sur chacune des lignes suivantes : le nombre
de \token{|} doit valoir \code{n-1} pour couvrir exactement la durée.

\begin{lstlisting}[caption={Fondu de 4 temps (1 ligne d'en-tête + 3 continuations).},numbers=none]
Ligne 1 :  fond_4_255_0_0_0_0_255
Ligne 2 :  |
Ligne 3 :  |
Ligne 4 :  |
\end{lstlisting}

\chapter{Dépannage}
\label{annexe:depannage}

\begin{description}[leftmargin=1.6em,style=nextline]
  \item[Le témoin DMX reste rouge / «{}DMX HS{}.]
    Le serveur ne répond pas : vérifiez qu'\Icare{} tourne encore et que
    l'adresse est bien \code{127.0.0.1:5000}.
  \item[La connexion échoue avec «{}No USB-DMX adapter found{}.]
    Aucun port COM ne contient «{}USB{}» dans sa description. Saisissez le port
    manuellement (par ex. \code{COM3}) dans le champ \emph{Port}. Vérifiez aussi
    que le pilote de l'adaptateur est installé.
  \item[Les projecteurs ne réagissent pas.]
    Contrôlez deux points : le \emph{canal de base} doit correspondre à celui
    programmé sur le projecteur, et la ligne DMX doit être correctement
    chaînée (adaptateur, projecteurs, bouchon de fin si nécessaire).
  \item[L'audio-réactif ne fonctionne pas.]
    Le navigateur doit avoir l'autorisation microphone. L'API
    \code{getUserMedia} n'existe que sur les pages servies en
    \code{localhost} ou en HTTPS : \Icare{} étant local, c'est le cas.
  \item[Le tempo mesuré est bizarre.]
    Les appuis trop espacés (plus de 3 secondes d'écart) sont ignorés.
    Utilisez \textbf{Reset} puis tapez de nouveau, plus régulièrement.
  \item[J'ai un ancien fichier de séquence.]
    Rien à faire : le format historique est migré automatiquement à l'ouverture.
\end{description}

\chapter{Glossaire}
\label{annexe:glossaire}

\begin{description}[leftmargin=1.6em,style=nextline]
  \item[Blueprints] Mécanisme Flask qui regroupe des routes par thème
        (\code{api\_live}, \code{api\_sequence}, \dots).
  \item[Canal de base] Premier canal occupé par un projecteur.
  \item[DMX512] Standard de pilotage de lumières : 512 canaux de 0 à 255 sur un
        bus RS-485 à 250 kbauds.
  \item[DSL] \emph{Domain Specific Language} : ici, le petit langage de chaînes
        qui décrit l'action d'un module sur un temps.
  \item[Enttec USB Pro] Adaptateur USB-DMX très répandu ; son protocole est
        celui implémenté par \file{pydmxIcare.py}.
  \item[Fixture] Un projecteur DMX.
  \item[Frame / image] L'ensemble des 512 valeurs envoyées sur la ligne.
  \item[HSV] Représentation d'une couleur par teinte, saturation et valeur ;
        pratique pour générer des arc-en-ciels.
  \item[Overlay] Effet appliqué par-dessus le mode courant (ici, le strobe).
  \item[RMS] \emph{Root Mean Square} : mesure d'énergie d'un signal sonore.
  \item[Séquence] Fichier JSON décrivant un show ligne par ligne.
  \item[Thread] Fil d'exécution. \Icare{} en utilise deux : la sortie DMX et la
        lecture Live.
  \item[Univers] Les 512 canaux d'une ligne DMX.
\end{description}


\end{document}
